Los resultados obtenidos responden directamente a los objetivos planteados en
la Sección~1.3. El procedimiento implementado --- formulación de dos índices
de Laporte y Nobert con generación iterativa de desigualdades RCI, delegando
cada modelo entero en el solver CBC --- demostró ser un método exacto
operativo: certificó el óptimo global en cuatro instancias estándar de
CVRPLIB, coincidiendo en todos los casos con las soluciones de referencia de
la biblioteca. Esta coincidencia exacta valida simultáneamente la formulación,
la convención de distancias adoptada y la correctitud de la separación de
cortes.

El hallazgo más relevante es la caracterización empírica del límite de
escalabilidad. La progresión de tamaños evaluada muestra tres regímenes
claramente diferenciados: las instancias pequeñas del Set E ($n \le 23$) se
resuelven en segundos; \texttt{A-n32-k5} ($n=32$) requiere cerca de diez
minutos pero cierra con certificado; y \texttt{A-n33-k5} ($n=33$) excede el
presupuesto sin certificar. Que un único nodo adicional multiplique el
esfuerzo por un orden de magnitud es la manifestación concreta del carácter
$\mathcal{NP}$-hard del problema discutido en la Sección~1.2, y de la
naturaleza exponencial de la familia de restricciones (5).

La principal limitación observada es estructural y no de implementación: el
esquema re-resuelve el modelo entero completo desde cero en cada iteración,
descartando todo el trabajo de búsqueda previo, y separa cortes únicamente
sobre soluciones enteras. La traza de \texttt{A-n32-k5} lo cuantifica: el
costo por resolución crece desde centésimas de segundo hasta superar los 16
segundos a medida que el modelo acumula cortes. Durante el desarrollo se
evaluaron mejoras de ejecución que no alteran el método: la separación
simultánea de múltiples cortes por iteración (componentes desconectadas,
rutas sobrecargadas y sub-bloques anidados) redujo las iteraciones de 333 a
99 y el tiempo en un 39\,\% sobre \texttt{A-n32-k5}, manteniendo el mismo
óptimo certificado, lo que confirma que la estrategia de separación gobierna
el costo total del esquema.

Estas observaciones conectan directamente con la progresión histórica revisada
en el Capítulo~2. La limitación de re-resolver desde cero es precisamente la
que resuelve \textit{Branch and Cut}, al separar cortes dentro de un único
árbol de búsqueda y sobre soluciones fraccionarias de la relajación lineal; y
la dificultad creciente de las instancias mayores motiva los esquemas
\textit{Branch, Cut and Price} que constituyen el estado del arte. Como
recomendaciones para trabajo futuro se identifican: (i) integrar la separación
de RCI dentro del árbol de búsqueda mediante \textit{callbacks} del solver, en
lugar del ciclo externo de re-optimización; (ii) separar cortes sobre
soluciones fraccionarias, fortaleciendo la relajación antes de ramificar; y
(iii) para instancias fuera del alcance de los métodos exactos, abordar el
problema mediante metaheurísticas, capaces de obtener soluciones de alta
calidad sin certificado de optimalidad, línea que se desarrollará en el
Informe Nro.~2.
